% Version 0.28.  The actual valued-point image-density reduction.
% Pinned source: 7cbf7a0c32fca1badf05931879a37c3f8a2b03b4.
\section{Valued points, nonsplit blocks, and the image-density reduction}
\label{sec:symplectic-image-density}

\status{IX1--IX8 prove the entire valued-point reduction to
\texttt{Fin0}, \texttt{NoEnd}, and automorphy with one common
complex embedding.  All of its local and global cohomological
inputs are supplied by AX, DX, BR, FP and LD; the full
\texttt{CohomologyHyp} record is unnecessary here.
IX9--IX10 also prove the characteristic-$\ell$ trace obstruction
and the characteristic-zero transport from an actual
\texttt{ProModAlt} witness.  These results do not construct
that witness.  The remaining seed, modular-support and
component-reaching assertions are recorded at the end with
their original quantifiers.  Consequently this chapter does
not assert that \texttt{SymplecticFourImageDense} has already
been supplied.}

\subsection{Fixed arithmetic data and the three exact incoming assertions}

Fix the original algebraic closure $\Omega/\mathbb Q$, an odd
prime $\ell$, the specified coefficient DVR
$O\subset\overline{\mathbb Q}_\ell$, and an actual good datum
$D$.  Write $k=O/\mathfrak m_O$, $E=\operatorname{Frac}O$,
$G=\operatorname{Gal}(\Omega/\mathbb Q)$, and retain every
chosen $G_p,I_p$ in $D$.  As in MX and PX, put
\[
 \theta=(\alpha\Psi,\Psi,\alpha,1),\quad
 r_0=\operatorname{diag}(\theta_i),\quad
 \lambda=(1,0,1,0),\quad \mu=(1,1,0,0),\quad
 w_{ij}=(\lambda_i-\lambda_j,\mu_i-\mu_j).
 \tag{IX.1}
\]
The cyclotomic exponents of $\theta$ on
$I_\ell\cap\ker\psi\cap\ker\omega$ are
$(e+1,e,1,0)$, with $e\ge2$.  We use exactly the alternating
form $J$ of (MX.3), and $X^*=J^{-1}X^tJ$.
Write $V_A=\langle e_0,e_2\rangle$ and
$V_B=\langle e_1,e_3\rangle$.  A block matrix preserves both
these summands.  They are isotropic and are paired perfectly.

Here a \emph{valued ambient field} is the source's
$\mathfrak a:\texttt{Amb}\,\ell$: an algebraically closed field
$K$ with a valuation ring $V$, compatible local maps
$a_{O'}:O'\to V$ for all permitted coefficient DVRs $O'$, and
the source's residue-field algebraicity condition.  A
representation, together with finitely many auxiliary matrices,
is \emph{small} if it is defined over a complete DVR $A\subset V$
with finite residue field, is continuous modulo each power of
$\mathfrak m_A$, and $K$ is algebraic over
$\operatorname{Frac}A$.  This is \texttt{Amb.Small}, including
all clauses of \texttt{IsDvr}; a merely bounded infinite set is
not being substituted for it.

The exact \texttt{KPair} consists of $r:G\to\operatorname{GSp}_4(V)$
and $u\in\operatorname{Sp}_4(V)$ such that:
the pair is small; $r$ reduces to $a_O r_0$;
its multiplier is $a_O\vartheta$; it satisfies
$D.\texttt{UnrT}$ and kills $I_p\cap\ker\psi$ for $p\in S_b$;
$u^{-1}r(G_\ell)u$ is upper triangular; its diagonal characters
reduce to the prescribed $\theta_i$ on $G_\ell$ and equal
$a_O\theta_i$ on $I_\ell$.
The exact \texttt{VPt} adds a witness consisting of an actual
$\texttt{IsPt}$ lift $\rho$ over some $R\in\mathcal C_O$ and a
local coefficient-compatible map $R\to V$ giving $r$; before
choosing $u$, smallness refers to $(r,1)$.

Set $\operatorname{Irr}(r)$ to mean that the $K$-span of the
matrices $r(g)$ is $M_4(K)$.  This is the source definition.
Set $\operatorname{BlockExt}(r)$ to mean that $r$ is block
triangular with its $A\leftarrow B$ block zero, its $B$-block
span equal to $M_2(K)$, and no matrix $Y\in\operatorname{Hom}(V_A,V_B)$
turns every $(1-Y)r(g)(1+Y)$ into a block diagonal matrix.
Put $\operatorname{Star}=\operatorname{Irr}\lor\operatorname{BlockExt}$.
The direction of the off-diagonal block is part of this definition.

Fix a field isomorphism
$\iota:\overline{\mathbb Q}_\ell\simeq\mathbb C$.
The three incoming assertions are exactly these:
\begin{enumerate}
\item $\operatorname{Con}(D)$: over every ambient field of
characteristic $\ell$, every actual \texttt{KPair} has
$\neg\operatorname{Star}(r_K)$.
\item $\operatorname{PurityAtPoints}(D,\iota)$: every
\texttt{KPair} with $\operatorname{Star}$, in every ambient
field identified with $\overline{\mathbb Q}_\ell$ by the
source's coefficient- and valuation-preserving
\texttt{IsZero} isomorphism, satisfies \texttt{PureW}.
\item $\operatorname{AutmAlt}(D,\iota)$: every such
characteristic-zero \texttt{KPair} with full matrix span
satisfies the literal \texttt{IsAut4Alt}$\,\iota$.
\end{enumerate}
In particular the second assertion concerns all these
\texttt{KPair}s, including the new nonsplit pair constructed
below; it is not just a purity statement for original
\texttt{EPt}s.  The source's \texttt{PureW} is the following
strict inequality, and is weaker than equality of absolute values:
there are an open $H\le G$ and $M\ne0$ such that, at every
prime $p\nmid M$, for every original place over $p$ and every
arithmetic Frobenius $\sigma\in H$ of degree $f>0$,
\[
 |\iota x|<p^f|\iota y|
 \quad\text{for all roots }x,y\text{ of }
                 \det(X-r(\sigma)),\quad y\ne0.
 \tag{IX.2}
\]
We keep \texttt{IsAut4Alt} as the literal \texttt{IsAutQA}
of the source, including its quadratic CM extensions,
regular algebraic cuspidal attachment and \texttt{WitA}
clauses.  A generic assertion of potential automorphy is
not substituted for it.

The additional classical facts used below are Noetherian
completion and Hensel's lemma, the weak Nullstellensatz,
properness of the symplectic flag variety, finite extensions
of complete discrete valuation fields, and Minkowski's theorem
that $\mathbb Q$ has no nontrivial finite extension unramified
at all finite primes.  The precise forms are recorded in the
classical-input registry.  For the geometric facts see Stacks
Project tags \href{https://stacks.math.columbia.edu/tag/04GH}{04GH},
\href{https://stacks.math.columbia.edu/tag/01WC}{01WC} and
\href{https://stacks.math.columbia.edu/tag/01KF}{01KF};
for the two number-theoretic facts see Milne,
\emph{Algebraic Number Theory}, Theorems 4.9 and 7.38--7.40,
\url{https://www.jmilne.org/math/CourseNotes/ANT.pdf}.
All finite local/global duality, finiteness and Euler formulas
are the already matched DX1 and LD3 inputs on the original
groups and finite coefficient modules.

% BEGIN IX1
\begin{proposition}[IX1: one universal ring for all actual points]
\label{claim:IX1}
There is a complete Noetherian local $O$-algebra $U$ with the
original residue field, an actual universal $\texttt{IsPt}$ lift
$\rho^u$, its origin $x:U\to O$, and a finite weighted
presentation
\[
       O[[z_1,\ldots,z_d]]\twoheadrightarrow U
 \tag{IX.3}
\]
with homogeneous kernel.  The variables are actual matrix-entry
differences $\rho^u(g)_{ij}-\theta_i(g)\delta_{ij}$ of weight
$w_{ij}$.  Every actual $\texttt{IsPt}$ lift over every
$R\in\mathcal C_O$ is induced by its unique $O$-algebra map
$U\to R$.
The full-weight-zero ring $B=U_0$ is complete Noetherian local
with residue field $k$.  If $B$ is finite over $O$, then
$D.\texttt{Fin0}$ holds.
\end{proposition}

\begin{proof}
DX2 and FP2 give the universal continuous matrix lift on
$D.Gam$.  MX1 identifies its pullbacks with precisely the
original global lifts satisfying the required inertia-killing
equations.  It remains to check that \texttt{IsPt}, including
its \texttt{OnComp} clauses, is a closed condition on these
actual lifts.

For each original local group, use the integral universal
ring $U_v$ of DX5--DX7 for $\operatorname{condL}$ or
$\operatorname{condP}$.  Let $x_v:U_v\to O$ be its origin,
$\phi_v:U_v\to E[[t_1,\ldots,t_{d_v}]]$ its
characteristic-zero chart, and $P_v=\ker\phi_v$.
The chart is a domain, $P_v$ is prime, and DX5 supplies
$a_v\notin\ker x_v$ with $a_vP_v=0$.
For every $f:U_v\to R$ we claim
\[
 \operatorname{OnComp}(\rho^u_v\otimes_fR,r_{0,v})
                  \quad\Longleftrightarrow\quad f(P_v)=0.
 \tag{IX.4}
\]
If $f$ kills $P_v$, the domain $U_v/P_v$, its induced origin,
its universal local representation and its map to $R$ are the
required \texttt{OnComp} witness.  Conversely a domain witness
$C$ gives, by local universality, a map $h:U_v\to C$ with
its composite to $O$ equal to $x_v$.  Hence $h(a_v)\ne0$.
The equality $a_vP_v=0$ and the fact that $C$ is a domain
force $h(P_v)=0$, and then $f(P_v)=0$.
This proves (IX.4) even for nonreduced target rings.

On an arbitrary lift first impose the local polynomial
equations, obtaining its quotient ring and its unique map
from $U_v$.  Impose the image of $P_v$ in that quotient.
The inverse image of this ideal gives exactly its
\texttt{OnComp} condition by (IX.4).  Add the multiplier,
inertia and normalization equations.  There are finitely
many local groups, and the ideal generated by all equations
is finitely generated and closed.  Its quotient is nonzero
because the actual diagonal lift $r_0$ satisfies them, by
MX1.  The quotient therefore represents the full actual
point functor, not only its characteristic-zero tangent
functor or one preselected family of points.

Conjugation by the formal diagonal torus preserves this
condition.  The \texttt{OnComp} clause is preserved directly:
for a domain witness $C$, replace it by
$C[[T_1,T_2]]$, conjugate by
$\operatorname{diag}((1+T_1)(1+T_2),1+T_2,1+T_1,1)$,
and send $T_i$ to the desired unit minus $1$ in the target.
The origin sends the $T_i$ to zero and still gives $r_0$.
The other clauses are polynomial conjugation identities.
Take finitely many universal entry differences spanning the
relative cotangent space.  FP1 makes them topological
generators, and the formal-torus coefficient argument of FP3
makes their presentation kernel homogeneous, with the stated
entry weights and actual origin.

The finite-monomial proof in PX1 applies to any finite set of
weights: a fixed weight has finitely many coordinatewise
minimal exponent vectors; every exponent of that weight is
one of these plus a weight-zero exponent.  The nonzero minimal
weight-zero exponents generate their monoid.  Therefore $B$
is a quotient of a formal power series ring over $O$ on
finitely many positive-order invariant monomials.  It has the
stated completeness, locality and residue field.

Finally any entry monomial of total weight zero in any actual
lift is the image of an element of $B$.  Its generated
$O$-algebra lies in the image of $B$ under the universal map.
If $B$ is finite, that image is a finite $O$-module, and any
submodule is finite because $O$ is Noetherian.  This is
\texttt{Fin0} with its full quantification over $R$ and $\rho$.
\end{proof}
% END IX1

% BEGIN IX2
\begin{proposition}[IX2: ordinary flags on every small valued point]
\label{claim:IX2}
Every actual $\operatorname{VPt}(D,\mathfrak a,r)$ admits
$u\in\operatorname{Sp}_4(V)$ making it an actual
$\operatorname{KPair}(D,\mathfrak a,r,u)$.
Every original $D.\texttt{EPt}$ is such a valued point in the
standard ambient $\overline{\mathbb Q}_\ell$ with its valuation
ring.  These assertions do not require
$\operatorname{CohomologyHyp}$ or any automorphy input.
\end{proposition}

\begin{proof}
Choose the original $\texttt{IsPt}$ witness and its local
domain witness $C$ at $\ell$.  The map $C\to O$ giving the
origin makes $O\hookrightarrow C$ injective.  Over
$F=\operatorname{Frac}C$ the four prescribed eigenvalues
$\theta_i(s_0)$ are distinct by AX6.  The unnormalized
projector equations in \texttt{Ord}, after these nonzero
differences are inverted, give their ordered stable full
flag; the multiplier pairs the $i$th and $(3-i)$th lines,
so it is a symplectic flag.  Its inertia quotient characters
are exactly $\theta_i$.

Take the schematic closure of this generic flag in the
symplectic flag variety over $C$.  The latter is projective:
embed a full flag in the product of the Grassmannians and
impose the closed incidence and isotropy equations.  Its
closed generic-flag closure is proper over $C$ and has image
all of $\operatorname{Spec}C$, since that image is closed and
contains the generic point.  All stability and fixed-inertia
equations hold on this closure: on each tautological
quotient they are regular equations vanishing at its generic
point.  After the given map $C\to V$ its generic fibre has a
$K$-point, because $K$ is algebraically closed.  Properness
extends that flag over the valuation ring $V$.
All its subquotients are free direct summands; successive
choice of a primitive vector and its symplectic dual gives
an adapted symplectic basis.  Let $u$ be its matrix.

The residual labels also need justification.  For $g\in G_\ell$
the $i$th generic quotient eigenvalue $b_i(g)$ belongs to
$F$: the ordered eigenspaces of $\rho(s_0)$ are defined over
$F$.  It is integral over $C$, because it is a root of the
monic characteristic polynomial of $\rho(g)$.
The finite domain $C[b_i(g)]$ is local because $C$ is
complete henselian.  Evaluating the projector expression
at the origin sends $b_i(g)$ to $\theta_i(g)$; this is
well-defined because its denominator is a product of the
nonzero scalars $\theta_i(s_0)-\theta_j(s_0)$.
Initially this evaluation takes values in $E$.  Its values
are integral over $O$, so they belong to $O$, which is
integrally closed.  The resulting map $C[b_i(g)]\to O$
is local: the inverse image of $\mathfrak m_O$ is a
maximal ideal above $\mathfrak m_C$, and the finite
henselian domain has only one such maximal ideal.
Thus $b_i(g)-\theta_i(g)$ belongs to the maximal ideal of
$C[b_i(g)]$.  Modulo $\mathfrak m_C$ it is nilpotent.
The determinant trick for this finite $C$-module supplies
a monic polynomial for it whose nonleading coefficients
lie in $\mathfrak m_C$.
The same polynomial equation holds on the generic-flag
closure and hence on the specialized quotient line.
In the valuation ring its root is in $\mathfrak m_V$.
This proves exactly the required residual character label
on every element of $G_\ell$, even when the characteristic
is $\ell$ and the four special eigenvalues coincide.

The remaining local conditions pass through the domain
witness by their polynomial equations, including identity
on $I_p\cap\ker\psi$ for $p\in S_b$.
Smallness is preserved: start with the DVR $A$ defining $r$;
the finitely many entries of $u$ and $u^{-1}$ are algebraic
over $\operatorname{Frac}A$.  Enlarge to their finite field
extension and its valuation ring $A'$.  It is a complete
DVR with finite residue field, and $u\in\operatorname{Sp}_4(A')$.
The induced maps are local, the entrywise adic continuity
of $r$ survives, and $K$ remains algebraic over its fraction
field.  This is the complete \texttt{Small} assertion.
The standard theorem for finite separable extensions is
Milne 7.38; in equal characteristic the purely inseparable
part is obtained by adjoining $\ell$-power roots of a
uniformizer.  Since the residue field is finite and perfect,
its valuation ring is again a finite complete DVR.

For an original $\texttt{EPt}$ with map $R\to
\overline{\mathbb Q}_\ell$, the boundedness argument of PX6
puts every element of $\mathfrak m_R$ in the open unit ball
and proves continuity.  A finite formal presentation shows
that its image lies in the finite extension of $E$ generated
by the variable values.  Its ring of integers gives the
required small local map and original residual congruence.
The identity ambient isomorphism fixes every coefficient
and its valuation, so satisfies the literal \texttt{IsZero}.
This also explains why the auxiliary DVR does not replace
the original datum or its coefficient maps.
\end{proof}
% END IX2

% BEGIN IX3
\begin{proposition}[IX3: an actual nonconstant characteristic-$\ell$ curve]
\label{claim:IX3}
If $D.\texttt{Fin0}$ fails, there are an ambient field of
characteristic $\ell$, an actual valued point $r$, and an
entry monomial $x$ of total weight zero such that
\[
                   x^{\ell^n}\ne x\qquad(n>0).
 \tag{IX.5}
\]
In particular the conclusion is a small actual point on the
original functor, not an arbitrary point of a larger special
fibre.
\end{proposition}

\begin{proof}
Use $U$ and $B=U_0$ from IX1.  Then $B$ is not finite over
$O$.  If $B/\mathfrak m_OB$ were Artinian, some
$\mathfrak m_B^N$ would lie in $\mathfrak m_OB$; the two
adic topologies would be equivalent, and successive
approximation from a basis modulo $\mathfrak m_O$ would
make $B$ a finite $O$-module.  Thus choose a prime
$P\supset\mathfrak m_OB$, maximal among those different
from $\mathfrak m_B$.  Such a maximal prime exists by
Noetherianity.  For $t\in\mathfrak m_B\setminus P$, the
domain $C=B/P$ has only its maximal ideal above $(t)$,
so $\mathfrak m_C^N\subset(t)$ for some $N$.
It is complete, contains $k$, and the continuous map
$k[[T]]\to C$, $T\mapsto t$, is injective: a nonzero series
is $T^a$ times a unit.  Successive approximation modulo
$t$ makes $C$ finite over $k[[T]]$.
Embed its fraction field in an algebraic closure
$K$ of $k((T))$ and take the extended valuation ring $V$.
This gives a local $y:B\to V$ with $y(t)\ne0$ and
$y(\mathfrak m_OB)=0$; its values lie in one finite extension
of $k((T))$.

To make the original ambient coefficient system, identify
the residue field of the union of the original $\ell$-adic
coefficient DVRs with $\overline{\mathbb F}_\ell\subset K$,
extending the specified embedding of $k$.  Reduction gives
all the maps $a_{O'}$.  They are local and compatible, and
every residue element comes from a finite residue extension,
hence from an unramified coefficient extension.  Therefore
the full \texttt{Amb} condition holds, not just the choice
of $K$ and $V$.

We now lift the point of the invariant ring.  Write
$U^{\mathrm{gr}}=\bigoplus_wU_w$.  Projection to weight zero
is $B$-linear and sends $1$ to $1$, so the fibre
$U^{\mathrm{gr}}\otimes_{B,y}K$ is nonzero.  The
finite-monomial argument of PX1 makes it a finite-type
$K$-algebra: it is generated by the finitely many $z_i$.
The weak Nullstellensatz gives a map
$\xi:U^{\mathrm{gr}}\to K$ extending $y$.
For its nonzero variable values put $a_i=v(\xi(z_i))$.
Whenever $n_i\ge0$ are integers, not all zero, and
$\sum n_iw_i=0$, locality of $y$ gives
\[
                         \sum n_i a_i>0.
 \tag{IX.6}
\]
The strict linear inequalities
$a_i+\langle h,w_i\rangle>0$ consequently have a rational
solution $h$.  Indeed otherwise separation from the positive
orthant gives nonzero rational $q_i\ge0$ with
$\sum q_iw_i=0$ and $\sum q_i a_i\le0$; clearing denominators
contradicts (IX.6).  This is the same finite rational
separation argument proved in PX5.
Scale by a torus element with valuation vector $h$.
It fixes $B$ and makes every nonzero variable value have
positive valuation.

These finitely many values, together with the finite field
of definition of $y(B)$, lie in one finite extension of
$k((T))$.  Let $A'$ be its complete valuation ring.
Power series can now be evaluated there.  This evaluation
agrees with $y$ on the full weight-zero formal ring: both
maps are continuous and agree on its finite monomial
generators.  For another weight, the finite-monomial
decomposition gives agreement with the scaled $\xi$.
Thus evaluation kills each homogeneous part of a defining
relation of $U$.  To kill the whole relation, sum the
finitely many complete weight components which have a
monomial of degree less than $N$; the remaining series has
order at least $N$.  Continuity and $N\to\infty$ finish the
argument.  We have obtained an actual local map $U\to A'$,
preserving $y$ on $B$, and hence a small actual valued point.

The maximal ideal of $B$ is generated by
$\mathfrak m_O$ and finitely many positive-order invariant
monomials in the entry differences.  Since $y(t)\ne0$,
at least one of these monomials has nonzero image in
$\mathfrak m_V$.  Every such invariant monomial is an
$O$-linear combination of entry monomials of total weight
zero: subtracting the origin affects only diagonal entries,
whose weights are zero.  If all the latter values were
fixed by some positive power of Frobenius, their sums and
products and the finite coefficient field would have the
same property after taking a common multiple of the powers.
But a nonzero element of $\mathfrak m_V$ cannot satisfy
$b^{\ell^n}=b$: the factor $1-b^{\ell^n-1}$ is a unit.
This contradiction proves (IX.5).
\end{proof}
% END IX3

% BEGIN IX4
\begin{proposition}[IX4: all reducible patterns and their constant invariants]
\label{claim:IX4}
For any actual small valued point, in either characteristic,
the alternatives of (PX.14) hold: full matrix span,
\texttt{DiagType}, or the two $2$-element blocks with all
cross-block entries zero in one direction and both $02,20$
entries present.  In the third case there is a block diagonal
\texttt{KPair} whose $A$-block span is $M_2(K)$.
In characteristic $\ell$, every weight-zero entry monomial
of a \texttt{DiagType} point satisfies
$x^{\ell^n}=x$ for some $n>0$.
\end{proposition}

\begin{proof}
The polynomial projector proof of PX6 takes place over the
actual complete coefficient ring before specialization.
It supplies all $e_{ii}$ in the $V$-span of $r(G)$, and
remains valid in characteristic $\ell$.  Consequently the
relation $i\rightsquigarrow j$ defined by a nonzero entry
is reflexive and transitive, is reversed by
$(i,j)\mapsto(3-j,3-i)$, and has full matrix span exactly
when all pairs occur.

We spell out the arithmetic exclusion in the new valued
setting.  A two-element strongly connected component
$\{s,u\}$ crossing the $\lambda$ blocks gives a
two-dimensional subquotient and an ordinary stable local
line by IX2.  In one of its coordinate charts the slope
$v$ lies in $V$, interchanging $s,u$ if necessary.
In the DVR defining smallness, choose a least-valuation
nonzero coefficient $b=r(g_1)_{su}$ and write
$r(g)_{su}=b c(g)$ with $c(g)\in V$.
The reductions of the $c(g)$ are locally constant and
not all zero: continuity modulo the valuation of $b$
plus one proves the first assertion.  The component
matrix identity gives the residual cocycle
\[
 z(g)=\frac{\overline{c(g)}}{\bar\theta_u(g)},
 \qquad \gamma=\bar\theta_s\bar\theta_u^{-1}.
 \tag{IX.7}
\]
It cannot be a coboundary.  A coboundary would say
$\overline{c(g)}=a(\bar\theta_s(g)-\bar\theta_u(g))$.
Apply the integral linear combination of $r(g)$ giving
$e_{ss}$: its off-diagonal entry says $0=a$, whereas some
$\overline{c(g)}$ is nonzero.

Outside $T$ the original unramified condition kills the
cocycle on the inertia kernel specified by \texttt{XVan}.
At $p\in S_b$, AX6 supplies $\sigma_p$ with
$\bar\chi_\ell(\sigma_p)=1$ and
$\gamma(\sigma_p)\ne1$.  Hence both the invariants of
$k(\gamma)$ and of its cyclotomic dual vanish; DX1 local
duality and the Euler formula for $p\ne\ell$ give $H^1=0$.
The same holds after any finite extension of the residue
field.  At $\ell$, use the coordinate chart in which the
stable line is generated by $(v,1)$, and write the
two-dimensional matrix as
$\bigl(\begin{smallmatrix}a_g&bc(g)\\ \beta_g&d_g\end{smallmatrix}\bigr)$.
Stability says
$bc(g)=(d_g-a_g+\beta_gv)v$.
At $g=\sigma_\ell$ the factor in parentheses is a unit:
$d_g-a_g$ is a unit by AX6, whereas
$\beta_g\in\mathfrak m_V$ and $v\in V$.
It follows first that $v/b\in V$.
Now divide every stable-line equation by $b$.
The extra term is $\beta_gv(v/b)$, whose reduction is
zero because the opposite off-diagonal entry $\beta_g$
is in $\mathfrak m_V$.  Hence
$\overline{c(g)}=(\bar\theta_u(g)-\bar\theta_s(g))
\overline{v/b}$ for $g\in G_\ell$.
Eliminating $\overline{v/b}$ gives exactly
\[
 \overline{c(g)}(\bar\theta_u(\sigma_\ell)
                         -\bar\theta_s(\sigma_\ell))
 =\overline{c(\sigma_\ell)}(\bar\theta_u(g)
                         -\bar\theta_s(g)).
 \tag{IX.8}
\]
The denominator on the left is nonzero by AX6, so the local
cocycle is a coboundary also at $\ell$.
The residue field of an ambient is algebraic over
$\mathbb F_\ell$; a locally constant cocycle on the compact
group has finite image.  Its values therefore lie in one
finite extension of $k$.  Applying a $k$-basis coordinate
to this cocycle and all its local primitives reduces it to
the exact residual vanishing in $D.\texttt{Good}$, supplied
by BR7.  Recombining coordinates makes (IX.7) a coboundary,
the desired contradiction.

The remaining finite graph argument is elementary.
Strong components are permuted by $i\mapsto3-i$.
A component of size three would meet its reversed image,
so would be reversal-stable, impossible for this
fixed-point-free involution.  A component of size four
gives full span.  A component of size two cannot cross
$\lambda$ by the preceding paragraph, so the two components
are $\{0,2\}$ and $\{1,3\}$.  Their quotient graph cannot
have edges in both directions.  If there is no component
of size two, there is no directed cycle between distinct
vertices, which is \texttt{DiagType}.

In the two-block case erase all cross-block entries.
Block projection is multiplicative on matrices triangular
in the same block direction and commutes with the
symplectic adjoint.  It therefore gives a representation
with the same multiplier, residual representation and
inertia-killing properties.  For the ordinary flag use
the following explicit specialization.  Capture the
original $r,u$ in its small DVR $A$.  Multiply its only
possible cross block by an indeterminate $t$.  The
resulting $r_t$ is a representation over $A[t]$ with the
same multiplier, and over $\operatorname{Frac}(A)(t)$
is conjugate to $r$ by a one-parameter diagonal torus.
Its ordered local diagonal characters are the original
elements of $A$, independent of $t$.
Take the closure of this generic flag in the projective
symplectic flag variety over $A[t]$, including the closed
equations for those fixed quotient characters.  Its
proper image contains the generic point and hence all
of $\operatorname{Spec}A[t]$.  At $t=0$ there is a flag
over $K$ with those same ordered characters, and
properness extends it over $V$.
Choose an integral adapted symplectic basis and capture
its finitely many entries in a finite DVR extension as
in IX2.  This proves every \texttt{KPair} clause for the
block projection.  The existing $e_{00},e_{22}$ and both
nonzero $02,20$ entries give all four $A$-block matrix
units, so that block has full span.

Finally suppose the characteristic is $\ell$ and the
pattern is \texttt{DiagType}.  Each diagonal entry
$d_i(g)=r(g)_{ii}$ is a character, since no off-diagonal
two-step return contributes to a diagonal product.
Put $U_1=\ker\bar r\cap\ker\psi$ and let $m=[G:U_1]$.
For inertia at $\ell$, $r(s^m)$ is conjugate to an upper
unipotent matrix; its fourth nilpotent power is zero, so
$r(s^m)^{\ell^2}=1$.  At $p\in S_b$ it is already $1$,
and away from $T$ the \texttt{UnrT} condition gives the
same conclusion.  Conjugacy of places gives
$d_i(s)^{m\ell^2}=1$ on every inertia group.
Reduce the powered character $d_i^{m\ell^2}$ modulo each power of the small
DVR's maximal ideal.  Its finite continuous image then
defines an extension of $\mathbb Q$ unramified at every
finite prime, hence trivial by Minkowski's theorem.
Separatedness gives $d_i(g)^{m\ell^2}=1$ for all $g$.

There is an integral linear functional $L$ on the weight
lattice which is positive on every present nonzero
off-diagonal weight.  One explicit choice is obtained
from $a_i=|\{j:i\rightsquigarrow j\}|$,
$h_i=a_i-a_{3-i}$ and
$L(x,y)=(h_0-h_1)x+(h_0-h_2)y$.
For a present strict edge its value is $h_i-h_j>0$.
A nonzero monomial of total weight zero can therefore
use only diagonal entries.  Their values have finite
order, so in characteristic $\ell$ they belong to a
finite field.  The monomial, or zero if one factor
vanishes, is fixed by a positive Frobenius power.
\end{proof}
% END IX4

% BEGIN IX5
\begin{theorem}[IX5: finite Selmer sizes produce a nonzero integral cocycle]
\label{claim:IX5}
Let $A$ be a complete DVR with finite residue field of
characteristic $\ell\ne2$, uniformizer $\pi$, and let $T_A$
be a finite free $A$-module with an adically continuous
action $\rho$ of the original $G$.
Fix a finite prime set $S$ containing $\ell$, an original
complex conjugation $c$, the original decomposition group
$H=G_\ell$, and $s\in H$ such that
\[
 \rho(c)=-1,\qquad \chi_\ell(s)=1,\qquad
                       \rho(s)-1\text{ is an automorphism of }T_A.
 \tag{IX.9}
\]
Let $P\ne0$ be a saturated $A$-submodule of $T_A$ stable
under $H$.  Let $K_0\le\ker\rho$ contain inertia outside
$S$, and put
\[
 N=\{g\in K_0:g\text{ belongs to an inertia group away from }\ell\}.
 \tag{IX.10}
\]
Then there are an adically continuous nonzero cocycle
$z:G\to T_A$ and $y\in T_A$ satisfying
\[
 \begin{gathered}
 z(gh)=z(g)+\rho(g)z(h),\qquad z(c)=0,
                  \qquad z(N)=0,\\
 z(h)-(\rho(h)y-y)\in P\qquad(h\in H).
 \end{gathered}
 \tag{IX.11}
\]
This holds in mixed and equal characteristic.  Its only
arithmetic input is the exact finite duality, finiteness
and Euler package already installed in DX1 and LD3.
\end{theorem}

\begin{proof}
Put $M_n=T_A/\pi^nT_A$ and
$P_n=(P+\pi^nT_A)/\pi^nT_A$.  Saturation over a DVR makes
$T_A/P$ free, so
$|P_n|=q^{n d}$, where $q=|A/\pi|$ and $d=\operatorname{rank}P>0$.
The finite module $M_n$ is killed by $\ell^n$, since
$\ell\in\mathfrak m_A$; it is therefore a legitimate
finite module for the installed $\mathbb Z/\ell^{n+1}$
cohomology interfaces, even in equal characteristic.
No freeness over that coefficient ring is asserted.

On $G_\ell$ choose the local condition
\[
 L_\ell=\operatorname{im}
                  \bigl(H^1(G_\ell,P_n)\to H^1(G_\ell,M_n)\bigr).
 \tag{IX.12}
\]
The map in (IX.12) is injective.  Indeed $\rho(s)-1$
is bijective on $M_n$ and on its stable finite subgroup
$P_n$; if a $P_n$-valued cocycle is the coboundary of
$m\in M_n$, evaluating at $s$ puts $m$ in $P_n$.
The local Euler formula at $\mathbb Q_\ell$ gives
\[
 |H^1(G_\ell,P_n)|
       =|P_n|\,|H^0(G_\ell,P_n)|\,|H^2(G_\ell,P_n)|
       \ \ge\ |P_n|.
 \tag{IX.13}
\]
At every other prime in $S$ use the unramified local
condition.  Its cardinality is that of $H^0(G_p,M_n)$:
both are the kernel or cokernel of Frobenius minus $1$
on the same finite inertia-invariant group.

The Greenberg--Wiles cardinality identity obtained from
the normalized finite Poitou--Tate pairing now gives
\[
                         |\operatorname{Sel}(M_n)|\ge |L_\ell|.
 \tag{IX.14}
\]
Here are all the cancellations in that formula.
The global invariants and real-place invariants of $M_n$
are zero because $c$ acts by $-1$ and $2$ is a unit.
The local invariants at $\ell$ are zero by $s-1$.
The global invariants of the cyclotomic dual are also
zero: a fixed functional evaluated on
$(\rho(s)-1)m$ is zero, since $\chi_\ell(s)=1$,
and these elements exhaust $M_n$.
At the other finite places the unramified factors
$|L_p|/|H^0(G_p,M_n)|$ equal $1$.
The only factor left besides $|L_\ell|$ is the positive
cardinality of the dual Selmer group.  All groups in
this argument are finite by DX1, with the same local
decomposition-group transport as LD3 and KB6.

Each Selmer class has a unique representative with
$z(c)=0$: subtract a coboundary, using
$\rho(c)-1=-2$.  This representative vanishes on $N$.
At $p\ne\ell$ a locally unramified class is a
coboundary on inertia, and a coboundary vanishes on
inertia intersected with $K_0\le\ker\rho$.
Outside $S$ it is already unramified.
Condition (IX.12) says exactly that some $y_n\in M_n$
makes its restriction differ from $\delta y_n$ by
$P_n$-valued elements.  Consequently the finite set
\[
 Z_n=\{z:G\to M_n:\ z\text{ satisfies the finite version of
                                      (IX.11) for some }y_n\}
 \tag{IX.15}
\]
satisfies $|Z_n|\ge q^{nd}$.  Its finiteness follows
also directly from its injection into the finite global
$H^1$: normalized coboundaries are zero.

We must produce compatible nonzero elements; mere size
growth would not suffice.  Call $z\in Z_n$ primitive
if its reduction in $M_1$ is nonzero.  Suppose for a
fixed $n\ge1$ that no element of $Z_n$ is primitive.
The same holds at every higher level by reduction.
For a nonprimitive $z\in Z_{m+1}$, divide its values
by $\pi$ to obtain a unique $z'\colon G\to M_m$.
This is a continuous cocycle, with the same zero
conditions.  To check its local condition, let $y$
be a witness for $z$.  At $s$, reduction modulo
$P+\pi T_A$ gives $(\rho(s)-1)y=0$.
The endomorphism $\rho(s)-1$ preserves $P$ and is
bijective on $P$ and $T_A/P$: its matrix is block
triangular for a splitting of the saturated submodule,
and its determinant is a unit.  Thus
$y=p+\pi y'$ with $p\in P$.
Replace $y$ by $\pi y'$ and divide the local equation
by $\pi$; saturation of $P$ gives the required condition
for $z'$ modulo $\pi^m$.
This division map $Z_{m+1}\to Z_m$ is injective whenever
all its elements are nonprimitive.
The resulting bounded sequence of cardinalities for
$m\ge n$ contradicts $|Z_m|\ge q^{md}$.
Hence each $Z_n$ contains a primitive element.

For each level take all pairs $(z_n,y_n)$ with $z_n$
primitive and $y_n$ a local witness.  They form nonempty
finite sets; reduction maps preserve both the equations
and primitivity for levels at least one.
The finite-branching tree of such pairs has paths of
every length, so has an infinite path: successively
choose a vertex with arbitrarily long continuations.
Completeness of the finite free module $T_A$ then gives
$z(g)=\lim z_n(g)$ and $y=\lim y_n$.
All the equations pass to the limit.
The submodule $P$ is closed since $T_A/P$ is finite
free and separated.  Adic continuity follows at each
finite level, and primitivity at level one gives
$z\ne0$.  This proves (IX.11) with no inverse-limit
surjectivity assumption.
\end{proof}
% END IX5

% BEGIN IX6
\begin{theorem}[IX6: the actual nonsplit block constructor]
\label{claim:IX6}
Let $(r,u)$ be any actual \texttt{KPair} for the original
good datum.  Suppose $r$ is block diagonal and its
$A$-block span over $K$ is $M_2(K)$.  Then there is another
actual \texttt{KPair} $(r',u')$ in the same ambient field
such that
\[
 \operatorname{BlockExt}(r'_K),\qquad
          \det(X-r'(g))=\det(X-r(g))\quad(g\in G).
 \tag{IX.16}
\]
No assertion that $r'$ is an original \texttt{EPt} is
needed or made.
\end{theorem}

\begin{proof}
Capture $r,u$ in their defining complete DVR $A$.
All identities hold there by injectivity into $V$.
Let
\[
 \mathcal U=\operatorname{Hom}(V_A,V_B)
                         \cap\mathfrak {sp}_4(A)
   =A(E_{10}-E_{32})\oplus AE_{12}\oplus AE_{30}.
 \tag{IX.17}
\]
It is free of rank three, its elements multiply to zero,
and it is stable under $\rho(g)Y=r(g)Yr(g)^{-1}$.
Put
\[
             P=\{Y\in\mathcal U:u^{-1}Yu
                                      \text{ is strictly upper triangular}\}.
 \tag{IX.18}
\]
It is saturated because its defining entries are linear
equations over a domain, and is stable under $G_\ell$
because upper triangular matrices preserve strict upper
triangularity by conjugation.

We check that $P\ne0$, rather than assuming it from a
dimension count.  Write $X=r(\sigma_\ell)$ and
$U=u^{-1}Xu$.  Residual cross-block separation makes
the $B$-block of $X-U_{00}$ invertible, so the first
column of $u$ is $v=(v_0,0,v_2,0)$.
Likewise $X|_{V_A}-U_{11}$ and $U_{00}-U_{11}$ are
invertible.  The first two column equations then give
the $A$-part of the second column as $a v$, for some
$a\in A$.  Put $w=u_1-av=(0,w_1,0,w_3)$.
The columns are independent, so $v,w\ne0$.
Isotropy gives $v_0w_3=v_2w_1$.
For $\varphi=(v_2,0,-v_0,0)$ the nonzero outer product
$Y=w\varphi$ lies in (IX.17).  It kills the first two
flag steps and has image in the second; hence
$u^{-1}Yu$ is strictly upper triangular.  This proves
$P\ne0$ over the original $A$.

Modulo $\mathfrak m_A$, the operator
$\rho(\sigma_\ell)-1$ on (IX.17) has the nonzero
diagonal eigenvalues
$\bar\theta_i(\sigma_\ell)/\bar\theta_j(\sigma_\ell)-1$
for the indicated cross-block positions.
Its determinant is therefore a unit, so it is an
automorphism over $A$.  AX6 also gives
$\chi_\ell(\sigma_\ell)=1$.
Complex conjugation acts on the two blocks as $-1,+1$:
each block of the involution $r(c)$ is congruent to that
scalar, and the other factor in $(r(c)-1)(r(c)+1)=0$
is a unit.  Thus
\[
                  r(c)=\operatorname{diag}(-1,1,-1,1),
                   \qquad \rho(c)|_{\mathcal U}=-1.
 \tag{IX.19}
\]
The original \texttt{UnrT} condition and AX6 make $r$
unramified outside $D.S$.  Its kernel lies in $\ker\rho$.
Every hypothesis of IX5 is now verified for $T_A=\mathcal U$,
$P$ as in (IX.18), and $K_0=\ker r$.

Let $z,y$ be the resulting nonzero cocycle and primitive.
Multiply both by $\pi$; they still satisfy (IX.11),
$z$ remains nonzero, and all entries of $z(g)$ lie in
$\mathfrak m_A$.  Define
\[
                       r'(g)=(1+z(g))r(g),
                       \qquad u'=(1-y)u.
 \tag{IX.20}
\]
Because products of two members of $\mathcal U$ vanish,
the cocycle identity is exactly the multiplicativity of
$r'$.  Also $(1+z)^*(1+z)=1$ and
$(1-y)^*(1-y)=1$; hence its multiplier and the symplectic
condition on $u'$ are unchanged.
The representation remains small and adically continuous,
with the same reduction.  If $g$ is in inertia away from
$\ell$ and $r(g)=1$, (IX.11) gives $z(g)=0$, so $r'(g)=1$.
This verifies the full original \texttt{UnrT} and all
$I_p\cap\ker\psi$ conditions, including their specified
kernel restrictions.

For $h\in G_\ell$ set
$W_h=z(h)-(r(h)yr(h)^{-1}-y)\in P$.
Using the same zero-product identities gives
\[
 (u')^{-1}r'(h)u'
           =(1+u^{-1}W_hu)(u^{-1}r(h)u).
 \tag{IX.21}
\]
This is upper triangular with exactly the old diagonal,
because the first added term is strictly upper triangular.
Thus the integral flag, residual labels and prescribed
inertia characters are all unchanged.  Every field of
the new \texttt{KPair} has been checked.

The matrix $r'$ is block triangular with the same two
diagonal blocks as $r$, proving the characteristic
polynomial identity.  The $B$-block has full span:
the symplectic multiplier identifies it with the
inverse transpose of the $A$-block, multiplied for
each $g$ by its nonzero similitude.  As $g$ ranges over
$G$, inversion and transpose preserve the full-span
condition, and these individual nonzero scalar factors
do not change the span.
If a cross-block $Y$ split the representation, expansion
of $(1-Y)r'(g)(1+Y)$ would give
$z(g)=Y-r(g)Yr(g)^{-1}$.
At $c$, (IX.19) and $z(c)=0$ give $2Y=0$, so $Y=0$;
then $z=0$, a contradiction.  The same argument works
after passage from $A$ to $K$, so the nonsplitting is
exactly the required \texttt{BlockExt} over $K$.
\end{proof}
% END IX6

% BEGIN IX7
\begin{proposition}[IX7: a block diagonal pair violates the exact weak purity bound]
\label{claim:IX7}
For every characteristic-zero ambient identification
satisfying \texttt{IsZero}, every block diagonal
\texttt{KPair} fails the literal \texttt{PureW} condition
(IX.2), for every fixed $\iota$.
\end{proposition}

\begin{proof}
Transport the pair to $\overline{\mathbb Q}_\ell$ using
the specified identification.  Smallness makes all its
entries continuous: an element in the $n$th power of
the defining DVR's maximal ideal has norm at most
$a^n$ for a fixed $a<1$.  Thus its two block determinants
$\delta_A,\delta_B$ are continuous characters.

The ordered ordinary flag gives
\[
 \delta_A|_{I_\ell}=\theta_0\theta_2,
                  \qquad
 \delta_B|_{I_\ell}=\theta_1\theta_3.
 \tag{IX.22}
\]
To justify the assignment of the two characters, use the
two blocks of $r(\sigma_\ell)$.  Their characteristic
polynomials are coprime over the valuation ring, since
their residual eigenvalues are disjoint by AX6.
The polynomial block projectors therefore preserve the
ordinary flag.  Their diagonal values on the flag
quotients are $1,0,1,0$ and $0,1,0,1$, as determined
by the original residual labels.  Multiplying the
appropriate quotient characters proves (IX.22).

Each of $\delta_A\chi_\ell^{-(e+2)}$ and
$\delta_B\chi_\ell^{-e}$ has finite image.
Here is the full global argument.  Choose an integer
$m>0$ killing the finite quotients attached to
$\bar r,\psi,\omega$.  On inertia at $\ell$, (IX.22)
makes the $m$th power of the relevant ratio trivial.
At $p\in S_b$ this follows from identity on
$I_p\cap\ker\psi$, and away from $T$ from
\texttt{UnrT}; the cyclotomic character is unramified
away from $\ell$.
Conjugacy of places gives the assertion at every finite
prime.  The characteristic-zero small DVR is finite
over $\mathbb Z_\ell$: it is complete, its residue field
is finite, and $\ell$ is a nonzero power of its
uniformizer times a unit, so successive approximation
from a basis modulo $\ell$ proves finiteness.
In this finite coefficient field defining the
pair, reduce the powered ratio modulo every maximal
ideal power.  Minkowski's theorem makes each of these
everywhere-unramified finite characters trivial.
Separatedness shows the powered ratio itself is $1$.
A continuous character with image in the finite set of
$m$th roots of unity has open kernel.
Hence on some common open subgroup $H_0$,
\[
                       \delta_A=\chi_\ell^{e+2},
                       \qquad \delta_B=\chi_\ell^e.
 \tag{IX.23}
\]

Suppose (IX.2) held with $H,M$.  Choose a prime
$p\nmid \ell M$.  An arithmetic Frobenius in the
chosen decomposition group at $p$ has some positive
power $\sigma$ in $H\cap H_0$, since that subgroup
has finite index.  If the power is $f>0$, then
$\chi_\ell(\sigma)=p^f$ and $\sigma$ is an arithmetic
Frobenius of degree $f$ at that same place.
Let $a_1,a_2$ and $b_1,b_2$ be the roots of the two
block characteristic polynomials.  They are all
nonzero because both blocks are invertible.
Using (IX.2) for $(a_1,b_1)$ and $(a_2,b_2)$ gives
\[
 |\iota\delta_A(\sigma)|
                   <p^{2f}|\iota\delta_B(\sigma)|.
 \tag{IX.24}
\]
But (IX.23) makes the two sides respectively
$p^{f(e+2)}$ and $p^{2f}p^{fe}$, which are equal.
This is a strict-inequality contradiction.
It uses exactly \texttt{PureW}, including its finite
exceptional set, open subgroup, all-place and
positive-Frobenius-degree quantifiers.
\end{proof}
% END IX7

% BEGIN IX8
\begin{theorem}[IX8: the complete actual image-density reduction]
\label{claim:IX8}
For the original good datum and one fixed $\iota$,
\[
 \begin{gathered}
 \operatorname{Con}(D),\quad
 \operatorname{PurityAtPoints}(D,\iota),\quad
 \operatorname{AutmAlt}(D,\iota)\\
 \Longrightarrow\quad
 D.\texttt{Fin0}\ \land\ D.\texttt{NoEnd}\ \land\
 D.\texttt{AllIrr}(\texttt{IsAut4Alt}\,\iota).
 \end{gathered}
 \tag{IX.25}
\]
All cohomological, valued-point, curve and block-extension
steps in this implication are proved, including the
coefficient identifications.  The three displayed arithmetic inputs
are not themselves asserted by it.
\end{theorem}

\begin{proof}
If \texttt{Fin0} failed, IX3 would give a characteristic-$\ell$
actual valued point with a weight-zero monomial satisfying
(IX.5).  IX2 gives its \texttt{KPair}.
Apply IX4.  Full matrix span contradicts
$\operatorname{Con}(D)$.  In \texttt{DiagType}, IX4 makes
the chosen monomial fixed by a positive Frobenius power,
contradicting (IX.5).  In the two-block case, IX4 first
gives a block diagonal pair with full $A$-block span;
IX6 turns it into a \texttt{BlockExt} pair.  This also
contradicts $\operatorname{Con}(D)$.
Thus \texttt{Fin0} holds on all original actual lifts.

To prove \texttt{NoEnd}, suppose an original algebraic
point were block diagonal with full $A$-block span.
IX2 places it in the standard ambient as an actual
\texttt{KPair}.  IX6 gives a new \texttt{BlockExt} pair
with exactly the same characteristic polynomial at
every group element.  The stated
$\operatorname{PurityAtPoints}$ applies to this new pair.
Since (IX.2) concerns only those characteristic
polynomials, it transfers to the old block diagonal
pair.  IX7 rules this out.  This is why the scope of
$\operatorname{PurityAtPoints}$ was retained literally.

Finally take any original algebraic point with full
matrix span.  IX2 makes it a \texttt{KPair} in the same
standard ambient, so $\operatorname{AutmAlt}(D,\iota)$
applies.  A field isomorphism carries a full matrix
span onto a full matrix span: map a finite spanning
linear combination entrywise and map its scalar
coefficients by the same isomorphism.
Thus no restriction on the actual point has been lost.
The isomorphism $\iota$ was fixed before this arbitrary
point was chosen.  It is the one common embedding in
the conclusion, not a point-dependent choice.
\end{proof}
% END IX8

\subsection{The finite-weight-fibre obstruction and the actual promodular witness}

The next two results unpack the exact source witness rather
than treating ``promodular'' as a classical theorem.  A
\texttt{ProModAlt} witness consists of a tower $T$, a
permitted coefficient extension $O\subset O'$, an actual
\texttt{Space} $X$ for $T$, an ideal $J\subset X.B$, a map
$y:X.A\to V$, and a global representation $r_1$ over $V$.
Here $X.B$ is complete Noetherian local with residue
$k'=O'/\mathfrak m_{O'}$, $X.B\hookrightarrow X.A$ is
injective, and $X.A$ is a finite-type $X.B$-algebra.
The universal framed pair $P$ has characteristic
polynomials $c_g(X)\in X.B[X]$ and local characters
$\psi_{vi}(s)\in X.B$.  Their residual values are the
specified four characters of $D$, with the tower's
permutation of labels.  The field $F=T.F$ is totally
real and abelian over $\mathbb Q$, $K_0=T.K$ is imaginary
quadratic, and $L=FK_0$.

The witness includes the full, unchanged
\texttt{T.Fit}, \texttt{T.Deep}, and
\texttt{PotentialModularityTower}$\,\iota\,X\,J$.
It also includes all of the following equations:
\[
 \begin{gathered}
 y|_{O'}=a_{O'},\qquad
 y(\mathfrak m_{X.B})\subset\mathfrak m_V,\qquad y(J)=0,\\
 r_K\simeq (r_1)_K,\qquad
 r_1|_{G_F}=P\otimes_yV,\qquad
 r_1(g)^*r_1(g)=a_O\vartheta(g),\\
 y(\psi_{vi}(s))=a_O\theta_i(s)
       \quad(s\in I_{F,v}),
 \end{gathered}
 \tag{IX.26}
\]
where the last line uses the original label convention of
\texttt{ProModAlt}; the residual characteristic-polynomial
formula uses the tower's possibly permuted labels.
The five clauses of the potential-modularity tower record are
a finite special weight fibre, local vanishing at the
specified characteristic-$\ell$ curve points, local
vanishing at actual seed points, weak purity for the
specified arithmetic points, and literal
\texttt{IsAut4Alt} for their irreducible global descents.
Only the first, fourth and fifth clauses are used below.
The unused clauses remain part of the actual source
record; no existence of such a record is inferred from
their names.

% BEGIN IX9
\begin{proposition}[IX9: finite weight modules, residual traces, and the full obstruction]
\label{claim:IX9}
The finite-fibre clause used in the obstruction has a
complete algebraic supplier.  Namely, on an actual
\texttt{Fit}, \texttt{Deep} tower retain a
\texttt{Link}, its universal ordinary determinant ring
$R$, its universal inertial-character algebra $\Lambda$,
an ideal $J_0\subset R$, and the actual
\texttt{IsPhi}$\,0$ map $\phi:R\to B=X.B$.
If $R/J_0$ is finite over $\Lambda$, then for
$J=J_0B$ there is $N\ge0$ such that
\[
 \mathfrak m_B^N\subset
 I:=J+\mathfrak m_{O'}B+
       (\psi_{vi}(s)-1:s\in I_{F,v}).
 \tag{IX.27}
\]
Consequently, every actual characteristic-$\ell$
\texttt{KPair} admitting a \texttt{ProModAlt} witness
has neither full matrix span nor \texttt{BlockExt}.
The conclusion uses no automorphic purity or
\texttt{CohomologyHyp} assumption.
\end{proposition}

\begin{proof}
\emph{Finite module to finite special fibre.}
Take a prime $\mathfrak q\supset I$ and put
$C=B/\mathfrak q$.  It is a complete local domain with
residue $k'$, and $\mathfrak m_{O'}C=0$.
An element of $C$ integral over $k'$ must belong to
$k'$: subtract its residue in $k'$; if the difference
were nonzero, it would be algebraic over a field and
therefore invertible, although it lies in the maximal
ideal.  Write $\pi:B\to C$.

The map $\Lambda\to R\xrightarrow{\phi}B\to C$
sends every universal inertial character to a scalar
in $k'$.  Indeed the exact \texttt{IsPhi}$\,0$
equation is
\[
 \phi(\psi^R_{vi}(s))
        =\kappa(s)\psi_{e(v),i}(s|_F),
 \tag{IX.28}
\]
with $\kappa(s)\in(O')^\times$; the last factor is $1$
in $C$ for inertial $s$.  The \texttt{Link} identifies
the places and inertia groups in this formula.
These characters therefore equal those from the
residue map $\Lambda\to k'\to C$.
Both maps satisfy the same inertial lifting conditions:
they are obtained from actual local characters, their
residues agree, and their prescribed product and
multiplier equations are preserved.
The uniqueness clause of \texttt{LamObj.univ} makes
the two maps equal.  Thus the entire image of
$\Lambda$ is contained in $k'$.
Since $R/J_0$ is finite over $\Lambda$, every element
of the image of $R$ in $C$ is integral over $k'$.

For $g\in G_F$, its square belongs to $G_L$, since
$L/F$ is quadratic.  The characteristic-polynomial
part of \texttt{IsPhi}$\,0$ identifies the polynomial
of the universal determinant at a lift of $g^2$ with
$c_{g^2}$ after multiplying its roots by the coefficient
unit $\kappa(g^2)$.  Consequently every coefficient
of $\pi(c_{g^2})$ is integral over $k'$; multiplication
by an $O'$-unit does not affect integrality.
It follows that every coefficient of $\pi(c_g)$
is integral over $k'$ as well.  Here is a proof that
does not assume that $B\hookrightarrow X.A$ remains
injective after quotienting.  For formal roots
$z_1,\ldots,z_4$, each $z_i$ satisfies the monic
polynomial $\prod_j(T^2-z_j^2)$.  Thus each elementary
symmetric polynomial in the $z_i$ is integral over
the ring of elementary symmetric polynomials in the
$z_j^2$.  The resulting monic polynomial identities
are universal identities for the coefficients of
the characteristic polynomials of a matrix and its
square.  Apply them to $P(g)$ in $X.A$ and descend
the identities to $B$ using the original injection
$B\hookrightarrow X.A$; only then apply $\pi$.
Transitivity of integrality proves the assertion.

Every element of a chosen decomposition group of
$G_F$ at $\ell$ lifts to $G_L$: the splitting clause
of \texttt{T.Ok}, with $\ell\in D.S$, puts that
decomposition group in the subgroup fixing $K_0$.
Equation (IX.28), now for arbitrary decomposition-group
elements, shows that each $\pi(\psi_{vi}(s))$ is
integral over $k'$.  All characteristic-polynomial
coefficients and all these local character values
therefore equal their residues in $k'\subset C$.
The exact \texttt{Space.genB} uniqueness clause says
that two maps $B\to C$ agreeing on these quantities
are equal.  Compare $\pi$ with the residue map
$B\to k'\to C$.  It follows that
$\pi(\mathfrak m_B)=0$, hence
$\mathfrak q=\mathfrak m_B$.
Every prime over $I$ is therefore the maximal ideal.
Since $B$ is Noetherian, some power of $\mathfrak m_B$
lies in $I$: take powers of finitely many generators
lying in $\sqrt I$ and expand a sufficiently large
power of their ideal.  This proves (IX.27).
It supplies exactly \texttt{TPmFin}'s conclusion;
the source's additional \texttt{UnramifiedInputs}
argument is not used in this algebraic proof.

\emph{Finite fibre to constant characteristic polynomials.}
Now take an actual \texttt{ProModAlt} witness in
characteristic $\ell$.  Its \texttt{fin} clause is
(IX.27), possibly with a larger ideal $J$.
Every coefficient map $a_{O'}$ kills
$\mathfrak m_{O'}$: its prime kernel contains the
nonzero element $\ell$, and a nonzero prime of a DVR
is its maximal ideal.  Equation (IX.26) kills $J$.
The \texttt{T.Ok.ineL} condition says that the original
residual representation is trivial on $I_{F,v}$;
therefore (IX.26) also kills each $\psi_{vi}(s)-1$.
The map $B\to V$ thus kills $\mathfrak m_B^N$.
Its kernel is prime, so it kills $\mathfrak m_B$ and
factors through $k'$.
The actual \texttt{Space.cp\_red} and the conjugacy
in (IX.26) now give
\[
 \det(X-r(g))=\prod_{i=0}^3
       (X-a_O\theta_i(g))\qquad(g\in G_F).
 \tag{IX.29}
\]
The tower permutation of the four labels does not
change this product.

For any $g\in G$, choose $a>0$ such that $g^a\in G_F$.
The four elements $a_O\theta_i(g)$ have finite order,
since they belong to the finite residue field of $O$.
Equation (IX.29) for $g^a$ implies that every eigenvalue
$z$ of $r(g)$ is algebraic over $\mathbb F_\ell$:
its $a$th power is one of the algebraic elements
$(a_O\theta_i(g))^a$.  Its trace is therefore algebraic
over $\mathbb F_\ell$.  The difference
\[
              \operatorname{tr}r(g)
                         -\sum_i a_O\theta_i(g)
 \tag{IX.30}
\]
lies in $\mathfrak m_V$ by the original
\texttt{KPair.red}.  A nonzero element algebraic over
a finite field has finite multiplicative order and
is a unit in every valuation ring containing it.
Thus (IX.30) is zero.  Since the $\theta_i$ are
characters, we obtain
\[
 \operatorname{tr}r(xgh)=\operatorname{tr}r(xhg)
                    \qquad(x,g,h\in G).
 \tag{IX.31}
\]

\emph{The trace obstruction for both kinds of star point.}
Suppose a matrix representation $S$ of dimension
$d>1$ spans $M_d(K)$, and an invertible $Q$ satisfies
\[
 \operatorname{tr}(QS(xgh))
                         =\operatorname{tr}(QS(xhg))
                    \qquad(x,g,h\in G).
 \tag{IX.32}
\]
Linear extension in $x$ makes
$\operatorname{tr}(QM[S(g),S(h)])=0$ for every
$M\in M_d(K)$.  Taking
$M=Q^{-1}e_{ji}$ proves that every matrix entry of
the commutator is zero.  Linear extension in $g,h$
would make the entire matrix algebra commutative,
contrary to its two elementary matrix units.
Taking $Q=1$ in (IX.31) rules out full matrix span.

For \texttt{BlockExt}, write the diagonal blocks
as $A(g),B(g)$ in the order $V_A,V_B$, with $B$
spanning $M_2(K)$.  The distinguished $g_0\in G_F$
from \texttt{T.Ok.dist} has four distinct eigenvalues
by (IX.29).  Put
\[
 f_A(X)=X^2-\operatorname{tr}(A(g_0))X+\det A(g_0),
                 \qquad Q=f_A(B(g_0)).
 \tag{IX.33}
\]
The two block characteristic polynomials are coprime.
Bezout's identity and Cayley--Hamilton make $Q$
invertible.  Cayley--Hamilton on the $A$ block also gives
\[
 \begin{split}
 \operatorname{tr}(QB(y))={}&\operatorname{tr}r(g_0^2y)
     -\operatorname{tr}(A(g_0))\operatorname{tr}r(g_0y)\\
     &+\det A(g_0)\operatorname{tr}r(y).
 \end{split}
 \tag{IX.34}
\]
Apply (IX.31) to the three terms.  It gives (IX.32)
for the full $B$-block span, another contradiction.
This last argument uses only the triangular pattern
and the full block span; it therefore rules out the
source's stronger nonsplit \texttt{BlockExt} assertion
with no extra premise.  This proves the
finite-weight-fibre-to-star-obstruction chain.
\end{proof}
% END IX9

% BEGIN IX10
\begin{theorem}[IX10: exact characteristic-zero transport and the common embedding]
\label{claim:IX10}
An actual \texttt{ProModAlt} witness, together with
the source's \texttt{IsZero} identification, gives
the literal \texttt{PureW} assertion for $r$.
If $r,u$ is also an actual \texttt{KPair} with full
matrix span, it gives the literal
\texttt{IsAut4Alt}$\,\iota$ assertion for $r$.
Consequently, for one fixed original good datum and
one fixed $\iota$, the assertion
\[
 \begin{gathered}
 \text{every actual \texttt{KPair} with \texttt{Star},
 in every \texttt{Amb},}\\
 \text{admits the literal \texttt{ProModAlt}$\,D\,\iota$
 witness}
 \end{gathered}
 \tag{IX.35}
\]
implies the entire conclusion of (IX.25), on the same
datum and with that one common $\iota$.
This theorem proves the reduction from (IX.35);
it does not assert (IX.35).
\end{theorem}

\begin{proof}
Let $f:V\to\overline{\mathbb Q}_\ell$ be the
valuation-preserving coefficient-compatible
identification, and compose $y$ with $f$.
It is an $O'$-algebra map, since
$f a_{O'}$ is the original inclusion of $O'$.
It kills $J$, and the transported $r_1$ restricts
to $P\otimes_{fy}\overline{\mathbb Q}_\ell$ on $G_F$.

The weight hypothesis of \texttt{PotentialModularityTower}
must be checked with its actual normalization.  On the
open subgroup
\[
             G_F\cap\ker\psi\cap\ker\omega
 \tag{IX.36}
\]
and on inertia at every place over $\ell$, the last
line of (IX.26) becomes
$f y(\psi_{vi}(s))=\chi_\ell(s)^{w_i}$,
where $(w_0,w_1,w_2,w_3)=(e+1,e,1,0)$.
The original character formula at the chosen place
extends to every such place by conjugacy of places;
one-dimensional characters are unchanged by
conjugation.  The source's \texttt{T.Wts} takes
\emph{squares} of these local characters.
Its parameter is therefore precisely
$(2(e+1),2e,2,0)$, not $(e+1,e,1,0)$.
The \texttt{pur} clause now supplies (IX.2) for
$f r_1$.  Conjugacy in (IX.26) makes all characteristic
polynomials of $fr$ and $fr_1$ identical.  Every
quantifier of \texttt{PureW}, including the same
open subgroup and excluded integer, transfers directly.

For automorphy there are two more hypotheses to prove.
First, smallness supplies continuity of $fr$: a
uniformizer of its defining DVR has image of norm
strictly less than $1$, and congruence modulo its
$n$th power bounds each matrix-entry difference by
that norm to the $n$th power.  Conjugacy gives
continuity of $fr_1$.  The multiplier equation of
(IX.26), with coefficient compatibility of $f$,
gives exactly $D.\texttt{Mult}(fr_1)$.

Second, full matrix span on $G$ must remain full
on $G_L$.  We give the whole argument.
Write $\rho=fr$ and $H=G_L$.  The extension $L/\mathbb Q$
is finite abelian, since $F/\mathbb Q$ is abelian and
$K_0/\mathbb Q$ is quadratic.  Thus $H$ is normal and
$G/H$ is finite abelian.
The \texttt{T.Ok.g3} clause supplies, for each
nonzero root difference, an element of $G_L$ on
which that residual root character is not $1$.
In the source's notation the root character satisfies
$\operatorname{rt}(w_{ij})\chi_j=\chi_i$;
this follows by substituting the two torus characters
and $\chi_0\chi_3=\chi_1\chi_2$.
After undoing the tower permutation, the four
$\bar\theta_i|_H$ are consequently distinct.

Full matrix span makes the four-dimensional
representation of $G$ simple.  Its restriction to
$H$ is semisimple: choose a simple $H$-submodule;
the sum of its finitely many $G/H$-translates is
a nonzero $G$-submodule and hence the whole space.
A sum of simple submodules is semisimple.
If this restriction were reducible, its commutant
$\operatorname{End}_H(\overline{\mathbb Q}_\ell^4)$
would have dimension greater than $1$.
The finite abelian quotient acts by conjugation
on this commutant.  In characteristic zero that
action diagonalizes.  Its fixed part is
$\operatorname{End}_G=\overline{\mathbb Q}_\ell$
by Schur's lemma, so there is a nontrivial character
$\eta$ of $G/H$ and a nonzero eigenvector $T$
intertwining $\rho$ with $\rho\otimes\eta$.
Its kernel and image are $G$-stable, so $T$ is invertible.
For $g$ with $\eta(g)\ne1$ and any $h\in H$, equality
of traces under this intertwiner gives
\[
                        \operatorname{tr}\rho(gh)=0.
 \tag{IX.37}
\]
The map $f$ is injective.  Thus (IX.37) is already
an equality for $r$ over $V$, which can be reduced
modulo $\mathfrak m_V$.  The original residual
condition gives
\[
       \sum_{i=0}^3\bar\theta_i(g)\bar\theta_i(h)=0
                                  \qquad(h\in H).
 \tag{IX.38}
\]
Distinct characters of any group into any field are
linearly independent: in a relation with the fewest
nonzero coefficients, translate by an element on
which two characters differ and subtract a suitable
scalar multiple to obtain a shorter relation.
The four characters here are distinct, whereas
every coefficient $\bar\theta_i(g)$ is a nonzero
unit.  Equation (IX.38) is impossible.
The restriction is therefore simple, and its matrix
span over the algebraically closed field is full.
For completeness, the latter implication is the
finite-dimensional density argument: a proper
matrix subalgebra acting simply has scalar
commutant, and the semisimple quotient acts through
a single full matrix block; its radical acts as
zero on the simple faithful module, so the algebra
itself is that full block.
This is the finite-dimensional density theorem over
an algebraically closed field; see Etingof et al.,
\emph{Introduction to Representation Theory},
Theorem~2.5 and Section~2.5,
\url{https://math.mit.edu/~etingof/replect.pdf}.

Conjugacy transfers full span on $G_L$ to $fr_1$.
All hypotheses of the actual \texttt{aut} clause
are now verified, so it gives literal
\texttt{IsAut4Alt}$\,\iota$ for $fr_1$.
This predicate is invariant under isomorphism of
the Galois representation.  In its defining
\texttt{IsAutQA}, continuity and irreducibility
are invariant under a fixed conjugation; for each
quadratic CM extension the same cuspidal
representation and regular algebraic weight remain,
and its attachment intertwiner is composed with
that conjugation.  The same operation transports
the intertwiners in \texttt{WitA}.
Thus it gives exactly the stated predicate for $fr$.

Finally (IX.35) and IX9 give
$\operatorname{Con}(D)$.
The first part of this theorem gives
$\operatorname{PurityAtPoints}(D,\iota)$ on all
star pairs, including the pairs created by IX6.
The second part gives
$\operatorname{AutmAlt}(D,\iota)$ on all full-span
pairs.  IX8 now proves (IX.25).
Every construction used the original $D$, coefficient
maps and local groups.  The auxiliary $O'$ occurs
only inside the literal witness and is compatible
with the ambient maps; it does not change the
original quantified endpoint.
The single $\iota$ is fixed before any ambient or
point is chosen, exactly as required.
\end{proof}
% END IX10

\subsection{The remaining arithmetic frontier, with source quantifiers}

The completed results in this chapter retain exactly
the hypotheses printed in each statement: the
geometric/cohomological constructions IX1--IX7 and
the algebraic reductions IX9--IX10.  IX8 is the
complete implication from the three explicitly
stated arithmetic predicates.  No one of those
predicates, and no actual witness (IX.35), has been
manufactured by assuming that a Lean theorem named
``published'' is an applicable classical theorem.

The next source producer of (IX.35) separates
the original \texttt{OfCase} data into
\texttt{Ctx2} and \texttt{Ctx3}, with the same
$\Omega,O,D$ and one choice of $\iota$.
Its remaining inputs are the following.
\begin{enumerate}
\item \textbf{Actual modular support.}
The literal \texttt{PMExAlt} says: for every tower
$T$ satisfying \texttt{Fit} and \texttt{Deep}, and
every $\iota$, there is a coefficient threshold
$O_1$ such that for every permitted $O'$
containing both $O$ and $O_1$, and for every actual
$\texttt{Space}\ X$ over $O'$, there is an ideal
$J\subset X.B$ satisfying all five clauses of
\texttt{PotentialModularityTower}.
The threshold precedes both $O'$ and $X$.
The source constructs $J$ as a sum of images of
two modular ideals under two actual
\texttt{Link.IsPhi} maps.  IX9 closes its
finite-weight-module-to-\texttt{fin} step completely.
It does not construct the modular ideals, their
finite modules, the two links, or the other four
clauses.  In particular finite cohomological
windows from TW do not supply a modular support.

\item \textbf{Seed points with the prescribed residual entries.}
\texttt{SeedIIAlt} quantifies over every original
\texttt{Ctx2} datum and every
\texttt{Ok}, \texttt{Two}, \texttt{Deep} tower,
then over every $\iota$; it requires a coefficient
extension and an actual continuous residual
\texttt{SeedPairAlt} in \texttt{InV}, with the
specified $-w_{\rm Al}$ alternative or both
$w_{\rm Be},-w_{\rm Th}$ entries.
For \texttt{SeedIVAlt}, fix an original
\texttt{Ctx3} datum and an \texttt{Ok}, \texttt{Four}
tower $T_0$.  For every $\iota$ and every ambient
field, one must choose a prime $w$ splitting in
$T_0$ and an actual residual pair $s$ for the
frame with the added prime $w$, supported in
$n_M$.  This choice is made before the following
universal quantifier: for every tower $T\ge T_0$
with \texttt{Ok} and \texttt{Dies}$\,\ell\,w\,T.F$,
there must be a coefficient extension and actual
\texttt{SeedPairAlt} $Q$ whose reduction lies over
that same $s$, is in \texttt{InA}, and has both
the $-w_{\rm Be}$ and $-w_{\rm Th}$ entries.
These are exact seed constructions, including
regular algebraic cuspidal attachment after the
specified twist and the parity and gap conditions;
they are not supplied by generic potential
automorphy or by a theorem about already
automorphic points.

\item \textbf{Reaching the support from the actual star point.}
The case-two and case-three producers require
their actual tower and space constructions,
entry-pattern reductions, characteristic-$\ell$
curve connections and component chains.
The resulting point must satisfy every equation
of (IX.26), especially $y(J)=0$, its original
inertial-character equations and conjugacy of the
global representation.  Existence of a nonempty
modular support does not give these assertions.
The case-three source route includes
\texttt{RedA}, \texttt{YCoin}, \texttt{FibY},
\texttt{Chain4}, \texttt{ChainDesc4} and
the literal fourth-chain reduction.  This chapter records their
exact producer anchors but does not promote them
to reconstructed human proofs.
\end{enumerate}

The original top-level source obtains these
arithmetic objects through \texttt{HypothesesH'}
and the unramified-input record, printed Hida and
adjoint-lifting constructions, and the rank-two
and rank-four seed producers.
\texttt{HypothesesH'} contains exact literal
automorphy, cohomological, base-change, descent,
semisimplicity and quaternionic assertions; it
has not been treated as one named standard
theorem.  Even a complete supplier of
\texttt{CohomologyHyp} would not by itself supply
the modular ideals or seed and component chains.
Conversely the entire lower reduction proved
here already uses only DX1, DX2, DX5--DX7,
AX6, BR7, FP2--FP3, MX1, PX1/PX5/PX6 and LD3.
It needs no still-open Bloch--Kato record.

The exact unclosed endpoint is therefore the
actual arithmetic realization (IX.35), uniformly
for the original \texttt{Good}/\texttt{OfCase}
data and with one $\iota$ before the point
quantifiers.  Once it is supplied, IX8--IX10
give literal \texttt{SymplecticFourImageDense};
the later MX/PX model and point comparisons can
then use its \texttt{Fin0}, \texttt{NoEnd} and
common-embedding automorphy on the very same
datum.  This is a closed downstream dependency
chain with an explicit arithmetic frontier,
not a claim that the full image-density
arithmetic input, RH, or GRH has been proved
in the human reconstruction.
